\section*{第 \ref{ch:b3:quantum-formalism} 章——量子力学的形式体系}

\begin{solution}{exo:b3:quantum-formalism:1}
(a) $(4 + 1)/5 = 1$，$(1 + 1)/2 = 1$。(b) $\braket\varphi\psi =
(2 - \iu)/\sqrt{10}$；$\braket\psi\varphi = (2 + \iu)/\sqrt{10}$：
二者互为共轭。(c) $|(2 - \iu)|^2/10 = 1/2$。(d) 解
$\braket\psi\chi = 0$：$\ket\chi = (\ket1 - 2\iu\ket2)/\sqrt5$。
\end{solution}

\begin{solution}{exo:b3:quantum-formalism:2}
第一个：厄米；$\pm1$ 对应 $(\ket1 \pm \ket2)/\sqrt2$。第二个：
厄米；$\pm1$ 对应 $(\ket1 \pm \iu\ket2)/\sqrt2$。第三个：不是
厄米（共轭转置不同）。第四个：厄米；
$\lambda^2 - 5\lambda + 4 = 0$ 给出 $1$ 与 $4$，对应的特征向量为
$\big({-(1 - \iu)}\ket1 + \ket2\big)/\sqrt3$ 与 $\big(\ket1 +
(1 + \iu)\ket2\big)/\sqrt3$。
\end{solution}

\begin{solution}{exo:b3:quantum-formalism:3}
(a) $\braket\beta\alpha = \cos\beta\cos\alpha + \sin\beta\sin\alpha
= \cos(\beta - \alpha)$：概率为 $\cos^2(\beta - \alpha)$。(b)
强度 $\propto$ 光子数：马吕斯定律。(c) 每个光子都是一次
独立试验：方差为 $N\mathcal P(1 - \mathcal P)$
——噪声本身即证明光子的存在。(d) $|\cos\beta + \iu\sin\beta|^2/2
= 1/2$ 对任意 $\beta$ 成立：圆偏振光不偏好任何
横向轴。
\end{solution}

\begin{solution}{exo:b3:quantum-formalism:4}
(a) $(\hat x\hat p - \hat p\hat x)f = -\iu\hbar(xf' - (xf)') =
\iu\hbar f$。(b) $[\hat x^2, \hat p] = 2\iu\hbar\hat x$；$[\hat x,
\hat p^2] = 2\iu\hbar\hat p$。(c) 加上并减去 $\hat B\hat A
\hat C$。(d) 由 (c) 用归纳法得 $\iu\hbar\,n\hat p^{n-1}$——
$[\hat x, \cdot]$ 作用为 $\iu\hbar\,\partial/\partial p$，
即与 $x$ 的泊松括号在量子层面的对应。
\end{solution}

\begin{solution}{exo:b3:quantum-formalism:5}
(a) $|\braket{45^\circ}{V}|^2 = \tfrac12$，塌缩后，
$|\braket{H}{45^\circ}|^2 = \tfrac12$：总概率为 $\tfrac14$。(b) 按
$60^\circ$、$30^\circ$ 的次序依次通过（三步，每步 $30^\circ$）：
$(\cos^230^\circ)^3 = 27/64 \approx 0.42$。(c) 每步
$\pi/2N$ 以 $\cos^2(\pi/2N)$ 的概率通过：存活概率为 $[\cos^2(\pi/2N)]^N
= 0.25$、$0.60$、$0.88$，对应 $N = 2, 5, 20$。(d) 当 $N \to \infty$
时存活概率趋于 $1$：许多次轻微测量能无损地把态
转过 $90^\circ$——测量被当作一个
转向盘。
\end{solution}

\begin{solution}{exo:b3:quantum-formalism:6}
(a) $[\hat A, \hat B] = \begin{pmatrix}0 & -2\\ 2 & 0\end{pmatrix}
\neq 0$：不相容。(b) $\hat B$ 作用于 $\ket1$ 时，确定得到 $+1$，
无需塌缩；随后 $\hat A$ 给出 $\pm1$，
概率各为 $\tfrac12$。(c) $\hat A \to +1$ 之后态为
$(\ket1 + \ket2)/\sqrt2$；$\hat B$ 将其塌缩为 $\ket1$ 或
$\ket2$（各 $\tfrac12$）；无论哪种情形，最后的 $\hat A$ 给出
$\pm1$ 的概率为 $\tfrac12$：以概率
$\tfrac12$ 出现矛盾。(d) 对易的可观测量共享本征态：中间的测量
不会造成扰动，重复测量的结果
必然一致。
\end{solution}

\begin{solution}{exo:b3:quantum-formalism:7}
(a) $\Delta x = \sigma$，$\Delta p = \hbar/2\sigma$（宽度为 $\sigma$ 的高斯函数
的傅里叶变换，其宽度为 $1/2\sigma$，
在 $k$ 空间中）：乘积恰为 $\hbar/2$。(b) 柯西–施瓦茨不等式取等号：
$\hat\beta\ket\psi \propto \hat\alpha\ket\psi$，比值为纯虚数
——对 $x, p$ 而言，该微分方程
仅有高斯解。(c) $\Delta p \ge \qty{5.3e-25}{kg.m/s}$：$E \sim
\Delta p^2/2m_{\text{e}} \approx \qty{1}{eV}$——原子是
电子伏特级的装置，因为它们是埃量级的盒子。(d)
$\Delta v \ge \qty{5e-27}{m/s}$：永远不会有任何可观测的影响。
\end{solution}

\begin{solution}{exo:b3:quantum-formalism:8}
(a) $[\hat H, \hat x] = -\iu\hbar\hat p/m$ 与 $[\hat H, \hat p] =
\iu\hbar V'(\hat x)$ 给出 $\dd\langle x\rangle/\dd t = \langle
p\rangle/m$，$\dd\langle p\rangle/\dd t = -\langle V'\rangle$。(b)
内积中的变量代换表明其为厄米的；
$\hat\Pi^2 = \mathbb 1$ 迫使本征值为 $\pm1$。(c) $V(-x) =
V(x)$ 使 $\hat H$ 对宇称不敏感；非简并本征态必须
因此同时也是 $\hat\Pi$ 的本征态：偶或奇。(d) 对称双势阱的
近简并双重态：一个偶态、一个奇态
——即氨分子的反转对，因隧道效应劈裂，
并以 \qty{24}{GHz} 的频率辐射。
\end{solution}

\begin{solution}{exo:b3:quantum-formalism:9}
(a) 两个态共享能量 $E$：只报告能量仍留下两种
可能。(b) $\hat S$ 交换标号：厄米，平方为单位算符，
与对称的 $\hat H$ 对易；在该能级内，其本征态为
$(\ket{1,2} \pm \ket{2,1})/\sqrt2$，本征值为
$\pm1$。(c) $(E, +)$ 与 $(E, -)$：唯一的标号。(d) 简并是一个不完整的地址；
造成简并的对称性同时
提供了缺失的那一位。
\end{solution}

\begin{solution}{exo:b3:quantum-formalism:10}
(a) $\Delta E\,\Delta A \ge \tfrac12|\langle[\hat H, \hat
A]\rangle| = \tfrac\hbar2|\dd\langle A\rangle/\dd t|$；两边同除即得。(b)
时间是理论的参数，而非算符：该关系式
限定了任何可测量量在给定能量涨落下能以\emph{多快}的速度演化。
(c) $\Delta E \sim \hbar/2\tau = \qty{3.3e-8}{eV}$：
自然线宽为若干兆赫兹。(d) 每
秒可观测的变化只要求 $\Delta E \gtrsim \qty{5e-35}{J}$——对于
宏观能量，毫无约束：钟表可以走动。
\end{solution}

\begin{solution}{exo:b3:quantum-formalism:11}
(a) $\|\hat\alpha\psi\|^2 = \braket{\psi}{\hat\alpha^2\psi} =
\Delta A^2$（利用 $\hat\alpha$ 的厄米性）。(b) $\braket{\hat\alpha
\psi}{\hat\beta\psi} = \tfrac12\langle\{\hat\alpha,
\hat\beta\}\rangle + \tfrac12\langle[\hat A, \hat B]\rangle$：
第一项为实数（厄米部分），第二项为纯虚数。
(c) 由 $|z|^2 \ge (\operatorname{Im}z)^2$ 即得该定理；取等号要求
向量成比例，\emph{且}反对易子的
平均值为零。(d) $\hat A$ 的本征态有 $\Delta A = 0$，从而迫使
$0 \ge \hbar/2$：对可归一化的态这是不可能的。平面波
只有作为不可归一化的理想化对象
位于希尔伯特空间之外，才能“达到”该下界。
\end{solution}

\begin{solution}{exo:b3:quantum-formalism:12}
(a) 在 $\omega t = \pi/2$，即 $t = T$。(b) 每次检测间，
态已旋转了 $\pi/2N$；测得处于 $\ket1$ 的
概率为 $\cos^2(\pi/2N)$，并\emph{坍缩回}$\ket1$；各次
检测相互独立，故取乘积。(c) $0$，$0.53$，$0.88$，$0.98$：
观测足够密集时，锅中的水永远不会沸腾。(d) 投影假设
（P5）：每次观测都将演化重置到
其起跑线。
\end{solution}

\begin{solution}{pb:b3:quantum-formalism:1}
\textbf{1.} 旋转把一组标准正交对映为
另一组标准正交对：$\braket{\nu_e}{\nu_\mu} = -\cos\theta\sin\theta +
\sin\theta\cos\theta = 0$。
\textbf{2.} P6：自由飞行由 $\hat H$ 生成，其
本征态仅通过相位自主演化——无论
产生过程选取的是哪组基。
\textbf{3.} $\ket{\psi(0)} = \ket{\nu_\mu} = -\sin\theta\ket{\nu_1}
+ \cos\theta\ket{\nu_2}$。
\textbf{4.} $\ket{\psi(t)} = -\sin\theta\,\eu^{-\iu E_1t/\hbar}
\ket{\nu_1} + \cos\theta\,\eu^{-\iu E_2t/\hbar}\ket{\nu_2}$。
\textbf{5.} 整体相位在每个 $|\braket\cdot\cdot|^2$ 中均被消去：
保留 $\Delta\phi = (E_2 - E_1)t/\hbar$。
\textbf{6.} $E_2 - E_1 = (m_2^2 - m_1^2)c^4/2E = \Delta m^2c^4/2E$。
\textbf{7.} $\braket{\nu_e}{\psi(t)} = \eu^{-\iu E_1t/\hbar}
\sin\theta\cos\theta\,(\eu^{-\iu\Delta\phi} - 1)$。
\textbf{8.} $|\cdots|^2 = \sin^2\theta\cos^2\theta\,|{\eu^{-\iu
\Delta\phi} - 1}|^2 = \sin^22\theta\,\sin^2(\Delta\phi/2)$。
\textbf{9.} 无混合或无质量劈裂时，便无振荡。任何
观测到的振荡都证明 $\theta \neq 0$ \emph{且} $m_1 \neq
m_2$——中微子有质量。
\textbf{10.} 代入 $\Delta\phi = \Delta m^2c^4L/2\hbar cE$；
取 $L_{\text{振荡}}$ 时辐角为 $\pi$。
\textbf{11.} 两个味概率是归一化态在标准正交基下分量的平方，
二者之和为 $1$——
演化的幺正性保持了范数。
\textbf{12.} 只有\emph{相对}相位是可观测的，它
包含 $E_2 - E_1 \propto m_2^2 - m_1^2$：绝对质量相互抵消。
\textbf{13.} 以所给单位计，$\dfrac{L}{4\hbar cE}$ 为：
$\num{e3}\,\unit{m}/(4 \times \num{1.973e-7}\,\unit{eV.m} \times
\num{e9}) = 1.27$ 每 $\unit{eV^2}$。
\textbf{14.} 下方充分发展：$1.27\,\Delta m^2 \times 12800
\gtrsim 1$，即 $\Delta m^2c^4 \gtrsim \num{6e-5}\,\unit{eV^2}$；
上方可忽略：$1.27\,\Delta m^2 \times 15 \ll 1$，即 $\ll
\num{5e-2}$：大致在 $10^{-4}$--$10^{-2}\,\unit{eV^2}$ 之间。
\textbf{15.} $L_{\text{振荡}} = \pi E/(1.27\,\Delta m^2c^4) \approx
\qty{990}{km}$（取 \qty{1}{GeV}）：\qty{15}{km} 处未见振荡，
\qty{12800}{km} 恰为十三个完整周期——与观测到的
图样一致。
\textbf{16.} 对多个长度周期及能量分布取平均，
$\langle\sin^2\rangle = \tfrac12$：压低因子为 $\tfrac12\sin^22
\theta$；测得的一半压低要求 $\sin^22\theta \approx 1$，
$\theta \approx 45^\circ$——自然选取了最大混合。
\textbf{17.} $m_2c^2 \approx \sqrt{\num{2.5e-3}} = \qty{0.05}{eV}$：
比电子轻千万倍——这是已知最轻的物质，
而目前无人知道其原因。
\textbf{18.} 大亚湾：$1.27 \times \num{2.5e-3} \times 1.5/0.004
\approx 1.2$——量级为一，已调好；KamLAND：$1.27 \times
\num{7.5e-5} \times 180/0.004 \approx 4.3$——量级为一，对太阳劈裂而言：
每条基线都是调到某一个
$\Delta m^2$ 的干涉仪。
\textbf{19.} 问题：太阳预言的 $\nu_e$ 只
到达了三分之一。解决：缺失的三分之二以其他
味的形式到达，$\nu_e$ 已转化成了这些味（SNO 统计了
总通量，发现太阳并无过错）。
\textbf{20.} 振荡要求 $\Delta m^2 \neq 0$：至少有一种
中微子有质量——这是超出原始标准模型的
第一项实验室物理学成果。
\textbf{21.} 在 $10^7$ 米的距离上保持相位相干，意味着
途中\emph{几乎没有}任何粒子发生相互作用：截面如此之小，
需要数千吨水才能捕获少量粒子——这也是
超级神冈探测器使用五万吨水的原因。
\textbf{22.} 味算符在味基中是对角的，
而味基并非能量基：$[\hat H, \text{味}] \neq 0$——
味仅仅不是自由飞行的守恒量，而
能量和动量是。
\textbf{23.} 两个劈裂相差三十倍：对于任一给定的
$L/E$，其中一种振荡是活跃的，另一种则要么冻结，要么被完全平均掉——
每个实验看到的都是一个有效两能级系统。
\textbf{24.} 产生基 $\neq$ 传播基（
即旋转 $\theta$）；P6 给出两个相位；玻恩定则将
相对相位转化为概率。
\textbf{25.} $\theta \approx 45^\circ$ 且 $\Delta m^2c^4 =
\qty{2.5e-3}{eV^2}$ 给出，对 GeV 能量的 μ 子中微子，振荡长度
约为 \qty{1000}{km}：两能级线性代数，经由地球得到验证，
而一次通量减半的消失现象获得了诺贝尔奖。
\end{solution}
